\section{Analysis}
\label{sec:analysis}

CIPHER is most valuable when failures are treated as structured signals rather
than as a flat collection of incorrect strings. The analysis should therefore
be read as a \emph{failure geometry}: which stage breaks first, how errors
change with difficulty, and which seeds should be preserved as regression
cases. Some subsections below remain partially prospective pending
full-pipeline results, but the intended analytical frame is already clear.

\subsection{Stage bottleneck analysis}

We decompose failures by stage to identify where the recovery policy breaks
down.

\paragraph{S1 (program inference).}
On simple levels, S1 often saturates first: regex matching achieves high parse
F1 on canonical instructions. The remaining errors are not random; they
concentrate in paraphrastic or compositional instructions where the system must
recover the latent operator sequence rather than match surface forms.

\paragraph{S2 (tool-chain execution).}
Execution failures are rare on short chains but increase with operation count,
primarily because the model must realize a valid \texttt{ffmpeg} chain rather
than merely name the right operations.

\paragraph{S3 (terminal-state extraction).}
S3 is \emph{not} the binding bottleneck for frontier VLMs once the correct
frame is provided. The true difficulty lies in the coupling between temporal
localization and execution: many systems that can read the key still fail to
arrive at the correct frame.

\paragraph{Near-miss errors are informative.}
Near-miss predictions should not be discarded as noise. They reveal whether a
model reached the correct terminal state but misread one character, localized
the wrong frame, or recovered the wrong but visually plausible distractor. In a
diagnostic benchmark, these distinctions matter more than a single binary EM
label.

\subsection{Character-level confusion in the near-miss regime}

Under calibrated extraction settings, models produce consistent substitutions
among visually similar glyphs rather than random strings. These errors are
useful because they identify a readable-but-fragile regime in which the task is
still discriminative.

\subsection{Video-temporal failure taxonomy}

The DAVIS video-temporal slice makes the failure geometry concrete. Among the
78 modern video-temporal tasks, 67 currently have no exact match from any
evaluated backend, but only one task is an all-missing failure across all
models. This is the desired hard-but-diagnostic regime: most failures still
contain usable intermediate evidence.

The dominant failure families are:
\begin{itemize}
  \item \textbf{near misses}: 21 tasks have no EM but reach maximum process
        partial credit at least 0.5, indicating that a model recovered part of
        the operation/frame/fragment structure but failed the final string;
  \item \textbf{frame-localization weakness}: 24 tasks have maximum
        frame-localization credit below 0.5, showing that models often read
        from non-target frames or from the entire clip rather than the two
        target frames;
  \item \textbf{over-reading}: 12 tasks contain predictions longer than the
        gold key by more than two characters, usually because the model absorbs
        distractors or neighboring yellow text into the answer; and
  \item \textbf{hard swap-and-reverse cases}: the highest-difficulty frontier
        is dominated by \texttt{swap\_reverse\_chars} seeds with 7--8 frames,
        high jitter, distractors, and non-central placement.
\end{itemize}

This taxonomy directly motivates the dense metric. A flat EM score would mark
all of these examples as identical failures. Process partial credit separates
``no evidence'', ``found the operation'', ``localized the wrong frame'', and
``read the right evidence but transformed it incorrectly'', which are different
failure modes for an agentic multimodal system.

\subsection{Canonical failure suite}

For benchmark-driven development, we recommend preserving a curated suite of
20--50 canonical failure seeds with associated traces. Each entry should store:
seed, level, background regime, gold operator chain, predicted operator chain,
execution trace, temporal localization decision, predicted key, and error
family. This turns CIPHER from a static benchmark into a replayable regression
environment for agentic systems.

The current video-temporal seedbank is the first implementation of this idea.
It stores not only the winning candidates, but also diverse non-winning
candidates selected by a utilization--diversity criterion. This matters because
the most useful benchmark samples are not always the hardest candidates by EM:
some are slightly easier but much more interpretable, and some expose specific
operation or temporal-localization weaknesses that would be lost if we kept
only the scalar best-of-run candidate.
